\scanpage{16}
\mspage{5}
$M,N,\sigma$ forment une 2-catégorie \deleted{2-équivalente à celle}\inserted{qui est un 2-ps-tors sous la catégorie de Picard} des champs de Picard \uline{stricts} d'inv. $M,N$ (et 0 donc) -- la ``différence'' de \deleted{deux} $\unsure{\mathscr P|P}$ et $\mathscr P'$ (épinglés par $M,N,\sigma$) est donc une \unsure{Cat|cat.} de Picard stricte d'inv. $M,N$ -- d'invariant $\in\operatorname{Ext}^2(S;M,N)$)

\noindent ----

\deleted{\unsure{§}} Il reste encore à voir à quelles conditions un triple $(M,N,\sigma)$ provient d'un $\mathscr P$. \deleted{(L'obstruction serait-elle dans $\operatorname{Ext}^3(S;M,N)$ ?)}
% slanted
\margin{cette condition est satisfaite ssi $\mathscr P$ (d'invariant $(M,N,\sigma)$) admet une deuxième \deleted{\unsure{contrainte de comm.}} \unsure{Comm.|comm.} qui en fasse un champ de Picard \uuline{strict}.}
(C'est OK si $\sigma$ est la restriction d'une application \inserted{bilin} antisymétrique $M\times M\longrightarrow N$ $\varphi(x,y)=-\varphi(y,x)$ -- c'est le cas si $M_2$ \deleted{\unsure{$M$ tel $M_2$ soit un module projectif sur}} $\unsure{)}$ est un \unsure{$\mathbb F_2$-}faisceau ``libre'', p.ex. si $M$ est ``libre'' sur $\mathbb Z$. Cela montre déjà que $M$ est quotient d'un $L$ tel que le composé $L\overset{p}{\unsure{\relbar|\longrightarrow}}M\xrightarrow{\sigma}{}_2N$ est réalisable par une cat. de Picard \inserted{$\mathscr P$} (déterminée \unsure{près|par|pour} un élément de $\operatorname{Ext}^2(S;L,N)$) La ``restriction'' de cette $\mathscr P$ à $(R,N)$ ($R=\operatorname{Ker}p$) définit une classe ds $\operatorname{Ext}^2(S;R,N)$ déterminée mod images d'éléments de $\operatorname{Ext}^{\unsure{2}}(S;L,N)$ donc son image \inserted{$\operatorname{ob}(\sigma)$} ds $\operatorname{Ext}^3(S;M,N)$ est can. déterminée, et elle est nulle ssi l'on
% archive sheet number: 15
